\documentclass[conference]{IEEEtran}
\IEEEoverridecommandlockouts
\usepackage{cite}
\usepackage{amsmath,amssymb,amsfonts}
\usepackage{algorithmic}
\usepackage{graphicx}
\usepackage{textcomp}
\usepackage{xcolor}

\begin{document}

\title{Scientific Writing for Research: Conference Article Framework\\
{\footnotesize \textsuperscript{*}Course Laboratory Milestone -- Week 3 Template Selection}
\thanks{Laboratory assignment for the Scientific Writing for Research course.}
}

\author{\IEEEauthorblockN{Flavian-Alexandru Manda}
\IEEEauthorblockA{\textit{Faculty of Automation and Computers} \\
\textit{Politehnica University of Timi\c{s}oara}\\
Timi\c{s}oara, Romania \\
student@upt.ro}
}

\maketitle

\begin{abstract}
This document demonstrates the integration of the standard IEEE conference template (IEEEtran) as required for the Scientific Writing for Research laboratory coursework. The manuscript adheres to the two-column format and prepares the formal structure following the IMRaD methodology (Introduction, Methodology, Results, and Discussion).
\end{abstract}

\begin{IEEEkeywords}
IEEEtran, scientific writing, research methodology, LaTeX template, conference paper
\end{IEEEkeywords}

\section{Introduction}
Scientific writing demands strict adherence to formal structures, rigorous citations, and reproducible methodologies \cite{ecarnot2015writing}.

\section{Methodology}
The methodology section outlines the research design, detailing theoretical foundations, experimental setups, or algorithmic architectures \cite{thuan2019construction}.

\section{Results}
Empirical or simulation metrics will be systematically reported within this section.

\section{Discussion}
The discussion contextualizes the findings within the broader academic literature, directly addressing research questions and limitations.

\section{Conclusion}
The final section summarizes key research contributions and outlines prospective avenues for future investigation.

\bibliographystyle{IEEEtran}
\bibliography{references}

\end{document}
